%==============================================================================
% Capítulo 3: Estado del Arte
%==============================================================================
\section{Estado del arte}\label{sec:estado-del-arte}

Este capítulo revisa la literatura académica y las soluciones comerciales disponibles en cuatro áreas relevantes para el presente trabajo: plataformas de mesa de ayuda, herramientas de orquestación de flujos, métodos de clasificación automática de tickets de soporte, y calibración de la confianza y cascadas de modelos. El propósito es situar la contribución de esta investigación en el contexto de los desarrollos previos y fundamentar la brecha que motiva el trabajo.

\noindent\textbf{Protocolo de revisión.} La revisión que sustenta este capítulo y los fundamentos del marco teórico se construyó con una búsqueda estructurada cuya fecha de corte es el 21 de octubre de 2026, sobre seis fuentes: Scopus, IEEE Xplore, ACM Digital Library, ACL Anthology, Google Scholar y arXiv. Las cadenas, formuladas en inglés, combinaron términos de dominio y de método: A1, \emph{("help desk" OR "service desk" OR "IT support") AND ("ticket classification" OR "incident classification") AND ("machine learning" OR "natural language processing" OR "large language model")}; A2, \emph{("large language model" OR "LLM") AND ("ticket" OR "incident") AND ("classification" OR "routing")}; A3, \emph{("Spanish" OR "español") AND ("large language model" OR "LLM") AND ("evaluation" OR "benchmark")}; A4, \emph{("Gemini" OR "Google Gemini") AND ("ticket" OR "incident" OR "help desk")}; y B1, \emph{("human-in-the-loop" OR "human in the loop" OR "human-AI interaction") AND ("machine learning" OR "large language model")}. Criterios de inclusión: publicación arbitrada 2023-2026 para el estado del arte moderno (A1-A4), idioma español o multilingüe y evaluación empírica; para los fundamentos de human-in-the-loop (B1) no se impuso restricción temporal. Criterios de exclusión: literatura no arbitrada, notas de prensa y material comercial sin respaldo metodológico. Una extensión posterior (grupo C), ejecutada el 25 de octubre de 2026 (su fecha de corte propia), agregó sobre las mismas seis bases dos cadenas: C1, {\shorthandoff{"}\emph{("calibration" OR "confidence estimation" OR "confidence elicitation" OR "uncertainty") AND ("large language model" OR "LLM" OR "language model" OR "neural network") AND ("classification" OR "few-shot" OR "verbalized" OR "overconfidence")}}; y C2, {\shorthandoff{"}\emph{("cascade" OR "cascading" OR "routing" OR "hybrid" OR "rule-based") AND ("large language model" OR "LLM" OR "natural language processing" OR "NLP") AND ("classification" OR "cost" OR "efficiency")}}. Se incluyeron publicaciones arbitradas, sin restricción temporal, sobre calibración o estimación de confianza de clasificadores y modelos de lenguaje (C1) o sobre cascadas y enrutamiento de modelos con criterio de costo y exactitud (C2), con evaluación empírica o método de referencia; se excluyeron los preprints sin versión arbitrada y el material no arbitrado. De las que cumplían se eligieron, hasta seis, las más próximas al escenario de este trabajo (un modelo consultado como servicio y una cascada de dos etapas). Cada consulta quedó registrada, con base, cadena, fecha y resultado, en la bitácora \texttt{docs/Tesis/Correcciones/bitacora-busqueda-c-110.md}, consistente con lo aquí declarado.

\subsection{Soluciones comerciales de mesa de ayuda}\label{subsec:soluciones-comerciales}

El relevamiento propio de la documentación pública de los proveedores, consultada el 19-10-2026, muestra que las soluciones comerciales de gestión de mesa de ayuda examinadas son integrales y cubren el ciclo completo de vida del ticket. Zendesk Suite ofrece funcionalidades de recepción multicanal, automatización de flujos básicos y, mediante un módulo adicional, capacidades de inteligencia artificial para sugerencia de respuestas y deflección de consultas frecuentes. Freshdesk propone una arquitectura comparable e incorpora el motor Freddy AI para clasificación automática y enrutamiento inteligente. Jira Service Management, parte del ecosistema de Atlassian, ofrece capacidades limitadas de automatización inteligente en su plan estándar; la edición Data Center permite el despliegue autoalojado para organizaciones con requisitos de soberanía sobre los datos. ServiceNow ITSM constituye la opción de referencia en el segmento corporativo, con un módulo Predictive Intelligence que aplica aprendizaje automático sobre el histórico de tickets de la organización.

El relevamiento propio de la documentación pública de los proveedores, consultada el 19-10-2026, también muestra que sus ofertas comerciales incorporan capacidades de inteligencia artificial descritas para organizaciones medianas, principalmente en mercados de habla inglesa. No obstante, estas plataformas presentan dos limitaciones relevantes para el contexto de la presente investigación. La primera es el modelo de costos basado en agente y mes, que escala linealmente con el tamaño del equipo de soporte. La segunda es que la mayoría de estas plataformas operan exclusivamente en modalidad de software como servicio (SaaS) con almacenamiento de datos en jurisdicciones extranjeras, lo que introduce consideraciones adicionales sobre transferencia internacional de datos personales bajo la Ley~25.326 \parencite{Ley25326_2000}.

En el segmento de código abierto, osTicket constituye una alternativa de despliegue autoalojado sin costo de licencia, aunque carece de capacidades nativas de clasificación automática mediante procesamiento de lenguaje natural moderno. Esta brecha funcional motiva la búsqueda de arquitecturas compositivas que integren componentes maduros de código abierto con servicios de inferencia bajo demanda.

\subsection{Plataformas de orquestación de flujos}\label{subsec:plataformas-orquestacion}

Dentro de las plataformas de orquestación de flujos relevadas para este trabajo (consulta 19-10-2026) se incluyen Zapier, Make (anteriormente Integromat), Microsoft Power Automate, Apache Airflow y N8N. Las primeras tres operan exclusivamente como SaaS, lo que las descarta como base para implementaciones donde el control sobre la infraestructura es un requisito explícito. Apache Airflow se orienta a flujos de trabajo programados de tipo procesamiento de datos por lotes, con un diseño centrado en grafos acíclicos dirigidos (DAG) definidos como código Python \parencite{Apache2024_airflow}; aunque potente, su paradigma de ejecución programada resulta inadecuado para un caso de uso reactivo y orientado a eventos como el de una mesa de ayuda.

N8N ofrece un modelo de construcción visual de flujos sobre nodos predefinidos, soporte nativo para webhooks y disparadores reactivos, ejecución basada en eventos y posibilidad de despliegue completamente autoalojado bajo Docker \parencite{n8n2024_docs}. La disponibilidad de su código fuente bajo licencia Sustainable Use facilita la auditoría de seguridad y la adaptación a requerimientos específicos sin generar dependencia exclusiva del proveedor. Estas características, verificadas en el relevamiento propio de la documentación pública de la plataforma, consultada el 19-10-2026, confluyen para posicionar a N8N como la opción más adecuada para el escenario propuesto, donde la soberanía sobre los datos y los flujos de procesamiento constituye un requisito no negociable.

\subsection{Trabajos académicos previos sobre clasificación de tickets}\label{subsec:trabajos-academicos}

La literatura académica sobre clasificación automática de tickets de soporte ha experimentado un crecimiento considerable durante los últimos años. La Tabla~\ref{tab:trabajos-academicos} sistematiza los trabajos más relevantes de esta subsección, organizados por enfoque metodológico y por orden cronológico.

\textcite{Paramesh2019_it_service_desk} comparan Naive Bayes multinomial, regresión logística, k vecinos más cercanos y SVM sobre tickets reales de una mesa de servicio de infraestructura de TI, y reportan buen desempeño de SVM en todas las muestras.

\textcite{Revina2020_it_tickets} abordan el mismo problema desde otra óptica: clasifican la \emph{complejidad} del ticket (baja, media o alta) y no su categoría. Con 30 y 60 tickets etiquetados, concluyen que los métodos simples y explicables, con pocos rasgos lingüísticos del dominio, son competitivos frente a alternativas más complejas. Excluyen además, por decisión de diseño y no como resultado medido, los \emph{word embeddings} preentrenados, con el argumento de que el vocabulario del dominio contiene numerosos términos ausentes de los corpus generales; esa tensión con el enfoque basado en LLM se retoma en la \autoref{subsec:comparacion-literatura}.

La literatura reciente sobre clasificación de tickets no respalda un desplazamiento generalizado hacia los LLM generativos: en las comparaciones directas relevadas, los modelos tipo BERT superan a los LLM en 172 clases en turco \parencite{paket2024itservice} y, ajustados con pocos datos, son más consistentes y menos costosos que los LLM \emph{zero-shot} \parencite{colavito2025benchmarking}.

\begin{table}[!ht]
\centering
\caption{Trabajos académicos previos sobre clasificación de tickets}
\label{tab:trabajos-academicos}
\small
\setlength{\tabcolsep}{4pt}
\begin{tabular}{@{}>{\raggedright\arraybackslash}p{2.0cm} >{\raggedright\arraybackslash}p{3.2cm} >{\raggedright\arraybackslash}p{3.5cm} >{\raggedright\arraybackslash}p{3.8cm} >{\raggedright\arraybackslash}p{2.1cm}@{}}
\toprule
\textbf{Trabajo} & \textbf{Enfoque} & \textbf{Corpus/idioma} & \textbf{Métrica reportada} & \textbf{Condición} \\
\midrule
\textcite{Paramesh2019_it_service_desk} & NB, LR, kNN y SVM & tickets reales de infraestructura TI; tamaño e idioma: no reportado en el alcance revisado & no reportado en el alcance revisado; cualitativo: SVM con buen desempeño & clasificación clásica supervisada \\
\textcite{Revina2020_it_tickets} & complejidad baja/media/alta; TF-IDF vs. rasgos lingüísticos; kNN, DT, NB, LR, SVM, semisupervisado & 28.243 y 4.684 tickets (30 y 60 etiquetados); inglés predominante & exactitud 1,000 (árbol de decisión, rasgos lingüísticos) frente a 0,433 como máximo con TF-IDF (primer conjunto) & comparación de representaciones y selección de rasgos \\
\textcite{zangari2023ticket} & revisión + BERT; inyección de jerarquía de etiquetas & 20.000 reclamos + 35.000 issues; idioma no reportado en el alcance revisado (texto inaccesible) & mejor embedding: $>$28\,\% en F1 sobre el estándar; ML-BERT: hasta 5,7\,\% en F1 y 5,4\,\% en exactitud (bugs) & clasificación multinivel, dos datasets \\
\textcite{paket2024itservice} & LLM (Titan, Llama, Mistral) vs. BERT; zero-shot, fine-tuning, RAG & turco, 172 clases & no reportado en el alcance revisado (texto cerrado); cualitativo: BERT supera a los LLM & clasificación multiclase de 172 clases \\
\textcite{liu2025tickit} (TickIt) & escalada de tickets con LLM; fine-tuning supervisado guiado por categoría & tickets de nube (ByteDance), 160 de prueba; idioma: no reportado en el alcance revisado & F1 0,862 (ajuste fino supervisado + CoT) frente a 0,801 como máximo con modelos pequeños & escalada online, dinámica y por relaciones \\
\textcite{colavito2025benchmarking} & benchmark de 22 LLM generativos + 2 BERT de base & dos datasets de issues de GitHub; tamaño: no reportado en el alcance revisado (texto cerrado) & no reportado en el alcance revisado (texto cerrado); cualitativo: BERT más consistente que los LLM generativos & zero-shot vs. fine-tuning \\
\bottomrule
\end{tabular}
\par\smallskip
{\footnotesize\textit{Fuente:} elaboración propia sobre las obras citadas (resumen y texto abierto; consulta 25-10-2026).\par}
\end{table}

La evaluación de modelos de lenguaje en escenarios de pocos ejemplos ha documentado que el rendimiento de clasificación depende fuertemente del idioma del corpus, con resultados que favorecen a los idiomas con más recursos \parencite{ahuja2023mega, xuan2025mmluprox}. Para el español y sus variantes regionales, \textcite{grandury2025laleaderboard} construyen una tabla comparativa de LLM específica para las variedades del español de España y Latinoamérica, mientras que \textcite{mayorrocher2026spanish} evalúan modelos con pruebas de español como lengua extranjera, \textcite{bernalbeltran2025mexicanspanish} comparan ajuste fino y aprendizaje en contexto sobre español mexicano, y \textcite{ramonferrer2026tripletotext} evalúan modelos de bajos recursos sobre un corpus en español. Esta línea evidencia que el español, y en particular sus variantes regionales como el rioplatense, ha recibido una atención notablemente menor en las evaluaciones sistemáticas publicadas. Cabe acotar, además, una brecha concreta en el alcance revisado: en las fuentes consultadas hasta la fecha de corte del protocolo de revisión de este capítulo (21 de octubre de 2026), con las cadenas de búsqueda allí declaradas, no se identificó artículo arbitrado 2023-2026 que evalue Google Gemini sobre tickets de soporte en español. Esta ausencia relativa al alcance revisado constituye una de las motivaciones académicas del presente trabajo, el cual aporta evidencia empírica sobre el desempeño del clasificador en español rioplatense aplicado a un dominio de mesa de ayuda de organización mediana.

\subsection{Calibración de la confianza y cascadas de modelos}\label{subsec:calibracion-cascadas}

\textcite{Guo2017_calibration} mostraron, en experimentos que incluyen clasificación de documentos, que las redes neuronales modernas tienden a la sobreconfianza y que esta se corrige a posteriori con un parámetro de temperatura. En la clasificación por instrucciones, \textcite{Zhao2021_calibrate_before_use} hallaron en GPT-3 y GPT-2 sesgos hacia ciertas etiquetas, que una calibración contextual de la distribución de salida reduce. Para modelos consultados solo por texto, \textcite{Tian2023_verbalized_confidence} hallaron, en preguntas y respuestas, mejor calibración de la confianza verbalizada que de las probabilidades condicionales, y \textcite{Xiong2024_confidence_elicitation} observaron sobreconfianza, valores concentrados en el extremo alto de la escala y poca capacidad para separar aciertos de errores.

En las cascadas, \textcite{Varshney2022_model_cascading} consultan modelos de capacidad creciente solo si la confianza de predicción del anterior no alcanza un umbral, con mejoras de eficiencia y exactitud en clasificación de comprensión del lenguaje natural; \textcite{Chen2024_frugalgpt} proponen una cascada de LLM de costo creciente con una función de puntaje aprendida y un umbral para aceptar o escalar. En las fuentes consultadas con C1 y C2 hasta el 25 de octubre de 2026 no se identificó trabajo arbitrado que evalúe una cascada de reglas determinísticas y un LLM con confianza verbalizada sobre tickets de soporte en español; la ausencia es relativa al alcance revisado.

\subsection{Análisis comparativo y brecha identificada}\label{subsec:analisis-comparativo}

La Tabla~\ref{tab:comparativa} compara las principales soluciones revisadas en tres dimensiones: modalidad de despliegue, clasificación automática y modelo de costo.

\begin{table}[htbp]
\centering
\caption{Comparación de soluciones para gestión automatizada de mesa de ayuda}
\label{tab:comparativa}
\small
\begin{tabular}{@{}>{\raggedright\arraybackslash}p{2.6cm} >{\raggedright\arraybackslash}p{2.8cm} >{\raggedright\arraybackslash}p{3.5cm} >{\raggedright\arraybackslash}p{3.8cm}@{}}
\toprule
\textbf{Solución} & \textbf{Despliegue} & \textbf{Clasificación automática} & \textbf{Modelo de costo} \\
\midrule
Zendesk Suite     & SaaS comercial      & Si, módulo de IA                & Alto, por agente/mes \\
Freshdesk         & SaaS comercial      & Si, motor Freddy AI             & Medio-alto, por agente \\
Jira Service Mgmt & SaaS o autoalojado  & Limitada en plan estándar       & Medio, por agente \\
ServiceNow ITSM   & SaaS comercial      & Si, Predictive Intelligence     & Muy alto, contrato anual \\
osTicket          & Código abierto      & No nativa                       & Sin costo \\
\bottomrule
\end{tabular}
\par\smallskip
{\footnotesize\textit{Fuente:} elaboración propia a partir de la documentación oficial de cada plataforma: Zendesk Suite (https://www.zendesk.com/service/help-desk-software/), Freshdesk (https://www.freshworks.com/freshdesk/), Jira Service Management (https://www.atlassian.com/software/jira/service-management), ServiceNow ITSM (https://www.servicenow.com/products/itsm.html) y osTicket (https://osticket.com/), consultadas el 19 de octubre de 2026.\par}
\end{table}

La síntesis comparativa evidencia que ninguna de las soluciones revisadas combina simultáneamente las tres propiedades que la tabla releva y que el contexto del trabajo demanda: despliegue autoalojado para resguardo de datos sensibles, costo bajo y predecible en el tiempo, y capacidad nativa de clasificación automática mediante NLP moderno. La brecha identificada, en el alcance revisado en este capítulo y hasta la fecha de corte declarada, se sintetiza en la ausencia de soluciones de bajo costo, autoalojadas, con clasificación basada en modelos de lenguaje grandes y validadas empíricamente sobre corpus en español rioplatense para el dominio específico de mesas de ayuda empresariales medianas.

La presente investigación se ubica precisamente en esa intersección y propone una arquitectura compositiva que reutiliza componentes maduros de código abierto y servicios de inferencia bajo demanda, manteniendo bajo control de la organización los datos sensibles y los flujos de orquestación, y aportando evidencia empírica reproducible sobre el desempeño del enfoque en un contexto lingüístico y organizacional escasamente documentado en la literatura, y sumando un canal telefónico con transcripción automática como rasgo propio, sin pretensión comparativa.

\newpage
